% TOP_PROBLEM: {"id": 3038, "title": "Long rainbow arithmetic progressions", "category_id": 2, "category": {"id": 2, "name": "combinatorics", "display_name": "Combinatorics", "description": "Counting problems, graph theory, discrete structures.", "slug": "combinatorics", "order_index": 2, "created_at": "2026-07-31T15:26:25.671Z"}, "status": "open"}
% FIRST_POSTED: null
% ENTRY_KIND: statement_only
% Imported from: batch14.tex; UnsolvedMath ID 3038
\documentclass[11pt]{article}
\usepackage[margin=25mm]{geometry}
\usepackage{fontspec}
\setmainfont{Latin Modern Roman}
\usepackage{amsmath,amssymb,mathtools}
\usepackage{unicode-math}
\setmathfont{Latin Modern Math}
\usepackage{xurl,hyperref}
\hypersetup{colorlinks=true,urlcolor=blue}
\setcounter{secnumdepth}{0}
\setlength{\parindent}{0pt}
\setlength{\parskip}{5pt}
\setlength{\emergencystretch}{3em}
\newcommand{\sref}[2]{\hyperlink{source:#1:#2}{[S#2]}}
\newcommand{\cataloguescope}[1]{\paragraph{Recorded target.}#1}
\begin{document}
% UnsolvedMath ID 3038; source catalogue code OPG-426
\setcounter{section}{3037}
\section{Long rainbow arithmetic progressions}\label{3038}
{\small\sffamily UnsolvedMath ID 3038 · Combinatorics}
\subsection{Definitions and mathematical statement}

\paragraph{Shared notation.}$\mathbb N=\{1,2,\ldots\}$, and $\mathbb Z,\mathbb Q,\mathbb R,\mathbb C$ denote the integers, rationals, reals and complex numbers. Any other conventions are specified locally.


For integers $k\geq3$, $t\geq1$ and $n\geq1$, an equinumerous $t$-coloring is a map $c:\{1,\ldots,tn\}\to\{1,\ldots,t\}$ with exactly $n$ integers of each color. A rainbow $k$-term arithmetic progression is a sequence $a,a+d,\ldots,a+(k-1)d$ lying in that interval, with integer $d\geq1$, on which all $k$ colors are distinct. Define $T_k$ to be the least $t$ such that every equinumerous $t$-coloring has a rainbow $k$-term progression for every $n\geq1$.
\paragraph{Statement recorded in the supplied source.}
Is $T_k=\Theta(k^2)$ as $k\to\infty$? Explicitly, do constants $c,C>0$ and $k_0$ exist such that $ck^2\leq T_k\leq Ck^2$ for every integer $k\geq k_0$? \sref{3038}{2}
\subsection{Short English statement}
Does the number of equal-size color classes needed to force a long rainbow arithmetic progression grow quadratically with its length?
\subsection{Sources}
\begin{description}
\item[\hypertarget{source:3038:1}{S1}] ULAM.AI and the UnsolvedMath Contributors, \emph{UnsolvedMath: A Curated Collection of Open Mathematics Problems}, record 3038 (OPG-426). CC BY 4.0. \url{https://huggingface.co/datasets/ulamai/UnsolvedMath}.
\item[\hypertarget{source:3038:2}{S2}] Jesse Geneson, \emph{The order of long rainbow arithmetic progressions}, arXiv:2607.15116 (2026), §1 and Theorem 1. Theorem 1 gives $T_k=\Theta(k^2\log k)$, refuting the recorded quadratic-order conjecture. \url{https://arxiv.org/abs/2607.15116}.
\end{description}
\label{end:3038}
\end{document}
